\chapter{Forward PV Computation}
\label{chap:wupy_pv}

The \textbf{Wu Python} track computes Ertel potential vorticity directly
on the regional $51\times87$ lat-lon grid using finite differences,
mirroring the Wu Fortran subroutine \texttt{PV\_CALC} in
\texttt{pvpialln\_94UV.f}.

\section{Ertel PV on Sigma Levels}

Ertel PV in isobaric-sigma coordinates is
\begin{equation}
  q = -g\left(
    \zeta_\theta\,\frac{\partial\theta}{\partial p}
    + \frac{\partial v}{\partial p}\frac{\partial\theta}{\partial x}
    - \frac{\partial u}{\partial p}\frac{\partial\theta}{\partial y}
  \right)
  \label{eq:ertel_sigma}
\end{equation}
where $\zeta_\theta = f + \zeta$ is the isentropic absolute vorticity and
all derivatives are on constant $\sigma$ surfaces.

The computation uses:
\begin{itemize}
  \item Vertical differencing: centred on interior levels, one-sided at
        $\sigma$-boundaries (levels 1 and NL)
  \item Horizontal differencing: centred 2nd-order on the $1.5^\circ$ grid
  \item Unit conversion: output in PVU ($10^{-6}\,$K\,m$^2$\,kg$^{-1}$\,s$^{-1}$)
\end{itemize}

\section{Balanced Streamfunction}

For each level $k$, the balanced streamfunction $\psi$ is obtained by
solving the 2-D Poisson equation
\begin{equation}
  \nabla^2\psi = \zeta = \frac{\partial v}{\partial x} - \frac{\partial u}{\partial y}
  \label{eq:psi_poisson}
\end{equation}
with Dirichlet boundary conditions $\psi = gH/f$ on the domain edges.
The solver is the red-black SOR described in \Cref{chap:fd}.

\section{Boundary Streamfunction from Wind Integration}

On the domain boundary (the 1-cell frame), $\psi$ is NOT set by the simple
geostrophic relation $\psi = gH/f$.  Instead, following Davis~(2022, eq.~2.40),
the wind field is integrated along the boundary:
\begin{equation}
  \psi(\mathbf{x}_\text{boundary}) = \psi_\text{NW}
  + \int_{\text{NW}\to\text{SW}} v\,dx
  + \int_{\text{SW}\to\text{SE}} u\,dy
  + \int_{\text{SE}\to\text{NE}} v\,dx
  + \int_{\text{NE}\to\text{NW}} u\,dy
  \label{eq:psi_boundary}
\end{equation}
starting from $\psi_\text{NW} = gH/f$ at the north-west corner and
integrating south $\to$ east $\to$ north $\to$ west, with a divergence
correction $d_\text{sum}$ to enforce periodicity.

This is implemented in \texttt{wu\_python/core/pv\_calc.py} as
\texttt{\_boundary\_psi\_from\_winds()}.  It reduces the $\psi$ discrepancy
vs.~the Fortran reference from $415\%$ to $\sim5\%$ (p99 relative error).

\section{Implementation}

Two main modules handle the forward PV calculation:

\begin{itemize}
  \item \texttt{wu\_python/core/pv\_calc.py} ---
        \begin{itemize}
          \item \texttt{compute\_ertel\_pv(U, V, TH, H, FC)}: full Ertel PV on
                all 10 $\sigma$-levels
          \item \texttt{invert\_vorticity\_balanced(VOR, H, U, V)}: Poisson-solve
                for $\psi$ with wind-integrated boundary conditions
          \item \texttt{\_boundary\_psi\_from\_winds(H, U, V)}: Davis eq.~2.40
        \end{itemize}
  \item \texttt{wu\_python/core/fd\_ops.py} ---
        \begin{itemize}
          \item \texttt{laplacian\_5pt(psi)}: 5-point discrete Laplacian
                (\ref{eq:laplacian_5pt})
          \item \texttt{relative\_vorticity(U, V)}: $\zeta = \partial_x v - \partial_y u$
        \end{itemize}
\end{itemize}

The pipeline step \texttt{05\_pass\_ab/run\_pass\_ab.py} computes both the
mean-state and event PV + $\psi$, saving them as \texttt{pass\_ab.npz} for
subsequent BALNC/BALP inversion.
